% sections/environnement.tex
% Phase 9: Environnement de travail (RAPP-02)

\section{Environnement de travail}
\label{sec:environnement}

Voici l'environnement mat\'eriel et logiciel utilis\'e pour le projet.

\begin{table}[H]
\centering
\small
\caption{Environnement de d\'eveloppement.}
\label{tab:environnement}
\begin{tabular}{p{4cm}lp{6cm}}
\toprule
\textbf{Outil} & \textbf{Version} & \textbf{R\^ole} \\
\midrule
Syst\`eme d'exploitation & Windows~11 & Plateforme de d\'eveloppement \\
Python & 3.10+ & Langage principal d'analyse \\
pip & 24.x & Gestionnaire de paquets \\
Git & 2.x & Contr\^ole de version \\
VS Code & 1.9x & \'Editeur de code \\
\TeX{} Live & 2024 & Distribution \LaTeX{} \\
pdf\LaTeX{} & (inclus) & Compilation du rapport \\
Biber & (inclus) & Backend bibliographique \\
latexmk & (inclus) & Automatisation compilation \\
Node.js & 20.x & Ex\'ecution du prototype web \\
npm & 10.x & Gestionnaire de paquets JavaScript \\
Vercel CLI & 51.x & D\'eploiement du prototype web \\
\bottomrule
\end{tabular}
\end{table}

Python a \'et\'e choisi pour ses biblioth\`eques scientifiques (NumPy, Pandas, NetworkX, Matplotlib). Les d\'ependances s'installent via \texttt{pip install -r requirements.txt}.

Le rapport est compil\'e avec pdf\LaTeX{}, plus simple que Lua\LaTeX{} ou Xe\LaTeX{} et suffisant pour un document en fran\c{c}ais avec l'encodage T1. \TeX{} Live fournit tous les paquets n\'ecessaires, et \texttt{latexmk} automatise les passes de compilation.

Le prototype web~\cite{cryptographexplorer_web} tourne sur Node.js et est b\^ati avec Vite. En local~: \texttt{npm install} puis \texttt{npm run dev}. Pour la mise en ligne~: \texttt{vercel deploy}. Les biblioth\`eques utilis\'ees c\^ot\'e front sont d\'ecrites dans la section~\ref{sec:bibliotheques}.
